\documentclass[12pt,a4paper]{article}
\usepackage[utf8]{inputenc}
\usepackage[T2A]{fontenc}
\usepackage[bulgarian]{babel}
\usepackage{geometry}
\usepackage{graphicx}
\usepackage{listings}
\usepackage{xcolor}
\usepackage{hyperref}
\usepackage{booktabs}
\usepackage{float}
\usepackage{tikz}
\usetikzlibrary{shapes,arrows,positioning}

\geometry{margin=2.5cm}

\lstset{
    basicstyle=\ttfamily\small,
    breaklines=true,
    frame=single,
    numbers=left,
    numberstyle=\tiny,
    keywordstyle=\color{blue},
    commentstyle=\color{green!60!black},
    stringstyle=\color{red}
}

\title{
    \vspace{2cm}
    \textbf{Курсов проект} \\
    \vspace{1cm}
    \Large{P2P Архитектури при Разпределено Динамично Балансиране на Паралелни Приложения}
}

\author{
    Изготвил: Георги Лазов \\
    Дисциплина: Разпределени ИТ Архитектури \\
    Специалност: Информационни системи
}

\date{Април 2026}

\begin{document}

\maketitle
\newpage

\tableofcontents
\newpage

\section{Въведение}

Паралелните и разпределени системи са ключови за решаването на изчислително интензивни задачи в съвременния свят. Един от основните предизвикателства при такива системи е ефективното разпределение на натоварването между процесорите/възлите, особено когато работното натоварване е нехомогенно или непредвидимо.

В централизираните архитектури (Master-Slave) един главен процес координира разпределението на задачите. Този подход е прост за имплементация, но създава ``тясно място'' (bottleneck) при мащабиране. Алтернативата е \textbf{Peer-to-Peer (P2P)} архитектура, където процесите комуникират директно помежду си без централен координатор.

Целта на този проект е да се изследва архитектурата на P2P паралелни приложения при разпределено динамично балансиране, като се анализират различни топологии на процесите - по-конкретно \textbf{верижна топология (chain)} и \textbf{пълен граф (all-to-all)}.

\section{Теоретична основа}

\subsection{Паралелни модели и SPMD}

При паралелното програмиране се различават два основни модела:

\begin{itemize}
    \item \textbf{SPMD (Single Program Multiple Data)} - всички процеси изпълняват една и съща програма, но върху различни данни. Това е най-разпространеният модел при MPI приложения.
    \item \textbf{MPMD (Multiple Program Multiple Data)} - различните процеси могат да изпълняват различни програми. Типичен пример са конвейрите за пакетна обработка.
\end{itemize}

В контекста на P2P динамичното балансиране работим основно с \textbf{SPMD модел}, където всички процеси са равнопоставени и могат да поемат произволна част от работата.

\subsection{Peer-to-Peer архитектури}

При P2P архитектурите няма централен координатор - всички възли (peers) са равнопоставени и комуникират директно помежду си.

\textbf{Предимства:}
\begin{itemize}
    \item Няма единична точка на отказ
    \item По-добра скалируемост
    \item Разпределено вземане на решения
\end{itemize}

\textbf{Недостатъци:}
\begin{itemize}
    \item По-сложна координация
    \item Възможно неоптимално разпределение
    \item Комуникационни разходи при пълен граф
\end{itemize}

\subsection{Топологии на процесите}

\subsubsection{Верижна топология (Chain/Ring)}

При верижната топология всеки процес комуникира само със съседите си (ляв и десен).

\textbf{Характеристики:}
\begin{itemize}
    \item Брой връзки: $n-1$ (за верига) или $n$ (за пръстен)
    \item Комуникационна сложност: $O(n)$ за глобална информация
    \item \textbf{Скалируем} - добавянето на нов процес изисква промяна само на съседните връзки
\end{itemize}

\subsubsection{Пълен граф (All-to-All)}

При пълния граф всеки процес може да комуникира директно с всеки друг.

\textbf{Характеристики:}
\begin{itemize}
    \item Брой връзки: $n(n-1)/2$
    \item Комуникационна сложност: $O(1)$ за директна комуникация
    \item \textbf{Не е скалируем} - комуникационните разходи растат квадратично
\end{itemize}

\subsection{Динамично балансиране на натоварването}

При динамичното балансиране работата се преразпределя по време на изпълнение. В P2P контекст използваме алгоритъма \textbf{Work Stealing}:

\begin{enumerate}
    \item Процес завършва задача
    \item Ако опашката му е празна, поисква работа от съсед
    \item Ако съседът има достатъчно работа, споделя половината
    \item Ако съседът също няма работа, продължава по веригата
\end{enumerate}

\section{Анализ на архитектурите}

\subsection{Централизирано vs. Разпределено балансиране}

\begin{table}[H]
\centering
\begin{tabular}{lcc}
\toprule
\textbf{Характеристика} & \textbf{Централизирано} & \textbf{Разпределено (P2P)} \\
\midrule
Скалируемост & Ограничена & По-добра \\
Сложност & Ниска & По-висока \\
Глобална информация & Да & Не (локална) \\
Отказоустойчивост & Ниска & По-висока \\
\bottomrule
\end{tabular}
\caption{Сравнение на подходите за балансиране}
\end{table}

\subsection{Сравнение на топологии}

\begin{table}[H]
\centering
\begin{tabular}{lcc}
\toprule
\textbf{Критерий} & \textbf{Верига} & \textbf{Пълен граф} \\
\midrule
Брой връзки & $O(n)$ & $O(n^2)$ \\
Латентност & $O(n)$ & $O(1)$ \\
Скалируемост & Висока & Ниска \\
Балансиране & Локално & Глобално \\
\bottomrule
\end{tabular}
\caption{Сравнение на топологиите}
\end{table}

\subsection{Хардуерни аспекти}

При облачни инфраструктури (напр. Amazon AWS) топологията на мрежата влияе пряко върху производителността:

\begin{itemize}
    \item \textbf{Placement Groups} - възлите в една група са с ниска латентност (<1ms)
    \item \textbf{Elastic Fabric Adapter} - специализиран интерфейс за HPC (до 100 Gbps)
    \item \textbf{NUMA архитектура} - трябва да се съобрази при многопроцесорни възли
\end{itemize}

\section{Проектиране}

\subsection{Функционално проектиране}

В рамките на проекта са реализирани три варианта:
\begin{enumerate}
    \item Статично циклично балансиране (baseline)
    \item Динамично P2P балансиране с верижна топология
    \item Динамично P2P балансиране с пълен граф
\end{enumerate}

\textbf{Проблем за тестване:} Рекурсивно изчисление на числата на Фибоначи - задача с нехомогенно натоварване.

\subsection{Архитектура на решението}

\textbf{Модел:} SPMD - всички процеси изпълняват един и същ код.

\textbf{Комуникационен протокол:}
\begin{itemize}
    \item \texttt{WORK\_REQUEST} - искане за работа
    \item \texttt{WORK\_RESPONSE} - отговор с работа
    \item \texttt{TERMINATE} - сигнал за край
\end{itemize}

\subsection{UML диаграми}

\subsubsection{Диаграма на последователността при верижна топология}

\begin{figure}[H]
\centering
\begin{tikzpicture}[node distance=2cm, auto]
    \tikzstyle{process} = [rectangle, draw, text width=1.5cm, text centered, minimum height=0.8cm]
    \tikzstyle{arrow} = [thick,->,>=stealth]

    \node[process] (P0) {P0};
    \node[process, right=of P0] (P1) {P1};
    \node[process, right=of P1] (P2) {P2};
    \node[process, right=of P2] (P3) {P3};

    \draw[arrow] (P1) -- node[above] {REQUEST} (P2);
    \draw[arrow, dashed] (P2) -- node[below] {RESPONSE} (P1);

    \draw[arrow] (P0) -- node[above, yshift=0.5cm] {REQUEST} (P1);
    \draw[arrow, dashed] (P1) -- node[below, yshift=-0.5cm] {NO\_WORK} (P0);
\end{tikzpicture}
\caption{Комуникация при верижна топология}
\end{figure}

\subsubsection{Диаграма на състоянията}

\begin{figure}[H]
\centering
\begin{tikzpicture}[node distance=2.5cm, auto]
    \tikzstyle{state} = [ellipse, draw, text width=2cm, text centered, minimum height=1cm]
    \tikzstyle{arrow} = [thick,->,>=stealth]

    \node[state] (init) {INIT};
    \node[state, right=of init] (working) {WORKING};
    \node[state, below=of working] (stealing) {STEALING};
    \node[state, right=of working] (term) {TERMINATED};

    \draw[arrow] (init) -- node[above] {start} (working);
    \draw[arrow] (working) edge[loop above] node[above] {task done} (working);
    \draw[arrow] (working) -- node[right] {empty queue} (stealing);
    \draw[arrow] (stealing) -- node[left] {got work} (working);
    \draw[arrow] (stealing) -- node[below] {no work} (term);
\end{tikzpicture}
\caption{Състояния на процес}
\end{figure}

\section{Имплементация}

\subsection{Технологичен избор}

\begin{itemize}
    \item \textbf{Език:} C
    \item \textbf{Паралелизъм:} MPI (Message Passing Interface)
    \item \textbf{Компилатор:} mpicc (OpenMPI)
\end{itemize}

\subsection{Програмен код}

Основните файлове са:
\begin{itemize}
    \item \texttt{static\_fine.c} - статично циклично балансиране
    \item \texttt{p2p\_chain\_fine.c} - P2P с верижна топология
    \item \texttt{p2p\_full\_fine.c} - P2P с пълен граф
    \item \texttt{p2p\_fib\_decomposition.c} - P2P с декомпозиция на Фибоначи
\end{itemize}

\section{Тестова среда}

\textbf{Хардуерна конфигурация:}
\begin{itemize}
    \item CPU: Apple M4 Pro (12 ядра: 8 performance + 4 efficiency)
    \item RAM: 48 GB
    \item L1 Cache: 192 KB (per P-core), 128 KB (per E-core)
    \item L2 Cache: 16 MB (споделен)
    \item OS: macOS Sequoia
\end{itemize}

\textbf{Параметри:}
\begin{itemize}
    \item Брой процеси (p): 1, 2, 4, 8
    \item Fibonacci стойност: 40 (за стандартни тестове), 45 (за декомпозиция)
    \item Повторения: 3
\end{itemize}

\section{Резултати}

\subsection{Статично балансиране (Baseline)}

При статичното балансиране задачите се разпределят предварително. Процес 0 получава 75\% от тежките задачи (fib(42-45)), което създава дисбаланс.

\begin{table}[H]
\centering
\begin{tabular}{ccccc}
\toprule
\textbf{p} & \textbf{$T_p$ (ms)} & \textbf{$S_p$} & \textbf{$E_p$} & \textbf{CPU баланс} \\
\midrule
1 & 652 & 1.00 & 1.00 & 100\% (1 ядро) \\
2 & 416 & 1.57 & 0.78 & P0:100\%, P1:45\% \\
4 & 447 & 1.46 & 0.36 & P0:100\%, P1-3:25\% \\
8 & 455 & 1.43 & 0.18 & P0:100\%, P1-7:15\% \\
\bottomrule
\end{tabular}
\caption{Резултати при статично балансиране (fib 40)}
\end{table}

\textbf{Наблюдение:} Процес 0 е натоварен 100\%, докато останалите процеси са idle повечето време. Ефективността пада драстично с увеличаване на броя процеси.

\subsection{P2P Пълен граф (Work Stealing)}

Work stealing позволява idle процеси да ``крадат'' задачи от натоварените.

\begin{table}[H]
\centering
\begin{tabular}{ccccc}
\toprule
\textbf{p} & \textbf{$T_p$ (ms)} & \textbf{$S_p$} & \textbf{$E_p$} & \textbf{CPU баланс} \\
\midrule
1 & 672 & 1.00 & 1.00 & 100\% (1 ядро) \\
2 & 380 & 1.77 & 0.88 & P0:95\%, P1:90\% \\
4 & 264 & 2.55 & 0.64 & Средно 70\% \\
8 & 261 & 2.57 & 0.32 & Средно 45\% \\
\bottomrule
\end{tabular}
\caption{Резултати при P2P пълен граф (fib 40)}
\end{table}

\textbf{Наблюдение:} Work stealing подобрява баланса - idle процеси ``крадат'' задачи от натоварените. Но ускорението все още е ограничено, защото всяка задача $fib(n)$ се изчислява изцяло от един процес.

\textbf{Проблем:} Най-тежката задача определя минималното време. Ако $fib(45)$ отнема 2.5s на един процес, никакво преразпределение не може да намали времето под 2.5s.

\subsection{P2P с декомпозиция на Фибоначи (Нов подход)}

За да постигнем истински паралелизъм, \textbf{декомпозираме} самите задачи:

\begin{itemize}
    \item Вместо един процес да изчислява цялото $fib(45)$
    \item Разбиваме: $fib(45) = fib(44) + fib(43)$
    \item $fib(44)$ и $fib(43)$ стават отделни задачи, разпределени на различни процеси
    \item Продължаваме докато $n \leq 30$ (праг за директно изчисление)
\end{itemize}

\begin{table}[H]
\centering
\begin{tabular}{cccccc}
\toprule
\textbf{p} & \textbf{$T_p$ (ms)} & \textbf{$S_p$} & \textbf{$E_p$} & \textbf{CPU баланс} & \textbf{Задачи} \\
\midrule
1 & 2028 & 1.00 & 1.00 & 100\% & 1597 \\
2 & 1034 & 1.96 & 0.98 & 100\% & 1597 \\
4 & 562 & 3.61 & 0.90 & 100\% & 1597 \\
8 & 359 & 5.65 & 0.71 & 99.8\% & 1597 \\
\bottomrule
\end{tabular}
\caption{Резултати при P2P с декомпозиция (fib 45, праг 30)}
\end{table}

\textbf{Наблюдение:} Всички процеси работят едновременно през цялото време. CPU балансът е близо до 100\%, което води до почти линейно ускорение.

\subsubsection{Алгоритъм на декомпозиция}

\begin{enumerate}
    \item \textbf{Начало:} Задача $fib(45)$ влиза в опашката на P0
    \item \textbf{Декомпозиция:} Ако $n > 30$:
    \begin{itemize}
        \item Създай частичен резултат (чака две деца)
        \item Генерирай задачи $fib(n-1)$ и $fib(n-2)$
        \item Едната задача остава локално, другата се изпраща на друг процес
    \end{itemize}
    \item \textbf{Директно изчисление:} Ако $n \leq 30$: изчисли рекурсивно
    \item \textbf{Комбиниране:} Когато и двете деца завършат $\rightarrow$ сумирай и изпрати на родителя
    \item \textbf{Край:} Когато коренната задача получи резултат
\end{enumerate}

\begin{figure}[H]
\centering
\begin{tikzpicture}[
    level distance=1.5cm,
    sibling distance=3cm,
    edge from parent/.style={draw,-latex},
    every node/.style={circle,draw,minimum size=8mm}
]
    \node[fill=blue!20] {45}
        child {node[fill=green!20] {44}
            child {node[fill=red!20] {43}}
            child {node[fill=yellow!20] {42}}
        }
        child {node[fill=red!20] {43}
            child {node[fill=yellow!20] {42}}
            child {node[fill=orange!20] {41}}
        };
\end{tikzpicture}
\caption{Декомпозиция на fib(45) - различни цветове = различни процеси}
\end{figure}

\section{Анализ на резултатите}

\subsection{Сравнение на трите подхода}

\begin{table}[H]
\centering
\begin{tabular}{lccc}
\toprule
\textbf{Метрика (8 процеса)} & \textbf{Статично} & \textbf{P2P Full} & \textbf{Декомпозиция} \\
\midrule
Ускорение $S_p$ & 1.43 & 2.57 & \textbf{5.65} \\
Ефективност $E_p$ & 0.18 & 0.32 & \textbf{0.71} \\
CPU баланс & 15-100\% & 45\% ср. & \textbf{100\%} \\
\bottomrule
\end{tabular}
\caption{Сравнение при 8 процеса}
\end{table}

\subsection{Основни наблюдения}

\begin{enumerate}
    \item \textbf{Статичното балансиране има таван $\approx 1.5x$}
    \begin{itemize}
        \item Причина: Процес 0 изчислява най-тежките задачи сам
        \item Останалите процеси idle $>75\%$ от времето
    \end{itemize}

    \item \textbf{Work Stealing подобрява до $\approx 2.5x$}
    \begin{itemize}
        \item Задачите се преразпределят динамично
        \item Но всяка $fib(n)$ все още е ``атомарна'' - един процес я изчислява
        \item Ограничението е най-тежката единична задача
    \end{itemize}

    \item \textbf{Декомпозицията постига $\approx 5.7x$ при 8 процеса}
    \begin{itemize}
        \item Задачите се \textbf{разбиват} на подзадачи
        \item Работата вътре в $fib(45)$ се разпределя между процесите
        \item 100\% използване на CPU ядрата
        \item Близко до линейно ускорение ($E_p = 0.71$)
    \end{itemize}
\end{enumerate}

\subsection{Защо декомпозицията е по-добра?}

При стандартните подходи, ако имаме задачи $fib(45), fib(35), fib(36), fib(37)$:
\begin{itemize}
    \item $fib(45)$ отнема $\approx 2.5s$ и се изчислява от един процес
    \item Останалите задачи завършват бързо ($< 0.5s$)
    \item Резултат: 3 процеса са idle докато P0 работи
\end{itemize}

При декомпозиция, $fib(45)$ се разбива на подзадачи:
\begin{itemize}
    \item $fib(45) = fib(44) + fib(43)$ - двете подзадачи отиват на различни процеси
    \item Процесът продължава рекурсивно докато $n \leq 30$
    \item Резултат: работата \textbf{вътре} в $fib(45)$ се разпределя между всички процеси
\end{itemize}

\subsection{Влияние на прага за декомпозиция}

Прагът определя кога да спрем декомпозицията и да изчислим директно:

\begin{table}[H]
\centering
\begin{tabular}{ccccc}
\toprule
\textbf{Праг} & \textbf{Задачи} & \textbf{Време (4p)} & \textbf{Overhead} \\
\midrule
20 & 17711 & 0.82 s & Висок (много малки задачи) \\
25 & 4181 & 0.58 s & Среден \\
30 & 1597 & 0.56 s & Оптимален \\
35 & 610 & 0.68 s & Нисък (едри задачи) \\
\bottomrule
\end{tabular}
\caption{Влияние на прага при fib(45), 4 процеса}
\end{table}

\textbf{Извод:} Праг 30 е оптимален - достатъчно фина грануларност за добър баланс, но не прекалено много комуникация.

\section{Заключение}

Проектът демонстрира три подхода за паралелно изчисление:

\begin{enumerate}
    \item \textbf{Статично балансиране} - просто, но неефективно при нехомогенно натоварване (ускорение $\approx 1.5x$)

    \item \textbf{P2P Work Stealing} - динамично преразпределя задачи, но ограничено от грануларността ($\approx 2.5x$)

    \item \textbf{Декомпозиция на задачите} - разбива работата на подзадачи, постига почти линейно ускорение ($\approx 5.7x$ при 8 процеса)
\end{enumerate}

\subsection{Ключови изводи}

\begin{itemize}
    \item \textbf{Грануларността е критична} - задачи, които не могат да се разбият, ограничават паралелизма
    \item \textbf{Рекурсивните алгоритми} като Фибоначи изискват специален подход (декомпозиция)
    \item \textbf{CPU балансът} е най-добрият индикатор за ефективност на паралелизма
    \item \textbf{Work Stealing} помага, но не решава фундаментални проблеми с грануларността
\end{itemize}

\subsection{Препоръки}

За постигане на линейно ускорение:
\begin{enumerate}
    \item Анализирай задачата - може ли да се декомпозира?
    \item Избери подходяща грануларност (не прекалено фина, не прекалено едра)
    \item Използвай work stealing за динамично балансиране
    \item Мониторирай CPU натоварването за идентифициране на дисбаланс
\end{enumerate}

\section*{Източници}

\begin{enumerate}
    \item Wilkinson, B., \& Allen, M. (2004). \textit{Parallel Programming}. Prentice Hall.
    \item Gropp, W., et al. (2014). \textit{Using MPI}. MIT Press.
    \item Kumar, V., et al. (2003). \textit{Introduction to Parallel Computing}. Addison-Wesley.
    \item Amazon Web Services. (2024). \textit{EC2 Placement Groups}. AWS Documentation.
    \item Blumofe, R. D., \& Leiserson, C. E. (1999). Scheduling by work stealing. \textit{JACM}, 46(5).
\end{enumerate}

\end{document}
